%% LyX 2.3.7 created this file.  For more info, see http://www.lyx.org/.
%% Do not edit unless you really know what you are doing.
\documentclass[english]{article}
\usepackage[T1]{fontenc}
\usepackage[latin9]{inputenc}
\usepackage{geometry}
\geometry{verbose,lmargin=3cm,rmargin=3cm,footskip=2cm}
\usepackage{amsmath}
\usepackage{babel}
\begin{document}
\title{Replication Algorithm for Variance Swap Pricing}
\author{Anurag Srivastava (anurag.srivastava@riskcare.com)}
\maketitle

\section{Summary}

\noindent

The document describes the details of pricing a swap based on a replication-portfolio
comprising of a log-contract and a stock or future position. The replication
determines the price based on the volatilities across the skew while
maintaining a constant dollar gamma for the portfolio. The main inputs
to the model are 
\begin{enumerate}
\item Stock prices or Futures Prices
\item Call options with strikes above and below a certain threshold close
to the initial stock price $S_{0}$
\item Put options with strikes above and below a certain threshold close
to the initial futures price $F_{0}$
\end{enumerate}
The stochastic volatility models are currently out of scope for the
current document. 

\section{Black Scholes Model}

\noindent

The basic idea for pricing is to construct a portfolio of European
options weighted as the inverse strike squared gives a constant exposure
to volatility. Notice that at each hedging interval, the PNL is given
as follows for gamma $\Gamma_{i}$, the return $R_{i}$ of the underlying,
the variance $\sigma^{2}$and the time-step size $\delta t$ 
\begin{align*}
\frac{1}{2}\sum_{i}\Gamma_{i}S_{i}^{2}(R_{i}^{2}-\sigma^{2}\delta t)
\end{align*}

Notice also that $\theta=-\frac{1}{2}\Gamma S^{2}\sigma^{2}$. A portfolio
proportional to $\frac{1}{S_{i}}$ would evidently keep a constant
second-derivative i.e. gamma $\Gamma$ ($\propto\frac{1}{S_{i}^{2}}$)
. More generally, we can say that the payoff $alnS+bS+c$ would have
a constant $\Gamma$. A variance swap can thus be replicated with
a log contract with options which are then delta-hedged. The following
expression summarise how a combination of puts and calls approximates
the payoff $alnS+bS+c$ (with price of underlying$S_{T}$at time $T$).

\begin{alignat}{1}
\int_{0}^{S_{0}}\frac{(K-S_{T})^{+}}{K^{2}}dK+\int_{S_{0}}^{\infty}\frac{(S_{T}-K)^{+}}{K^{2}}dK & =\frac{S_{T}}{S_{0}}-ln(\frac{S_{T}}{S_{0}})\label{eq:main_replication}
\end{alignat}

The replication can also be derived in terms of Stock prices using
Black-Scholes assumptions (see Section \ref{subsec:Stocks-based-replication}).

\subsection{Futures-based replication}

We now discuss how the replication summarised with Equation \ref{eq:main_replication}
can be performed using futures prices $F_{t}$ (at time $t$) and
$F_{0}$. As discussed in \cite{allenJPMvarianceswaps2006,carrmadanvoltrading1998},
the average variance over time $\frac{1}{T}\int_{0}^{T}\sigma_{t}dt$
can be written as 

\begin{align}
\frac{1}{T}\int_{0}^{T}\sigma_{t}dt & =\frac{2}{T}(log(\frac{F_{0}}{F_{t}})+\frac{F_{t}}{F_{0}}-1)-\frac{2}{T}\int_{0}^{T}(\frac{1}{F_{0}}-\frac{1}{F_{t}})dF_{t}\label{eq:replication_with_futures_main}
\end{align}

The first-term can be replicated by taking a static position in a
contract that pays $\frac{2}{T}(log(\frac{F_{0}}{F_{t}})+\frac{F_{t}}{F_{0}}-1)$
and the second-term can be replicated by a continuously rebalanced
futures position. As detailed in \cite{carrmadanvoltrading1998},
this can be performed using calls and puts with a position $\frac{1}{K^{2}}$
using the following expression

\begin{align*}
log(\frac{F_{0}}{F_{T}})+\frac{F_{T}}{F_{0}}-1) & =\int_{0}^{F_{0}}\frac{(K-F_{T})^{+}}{K^{2}}dK+\int_{F_{0}}^{\infty}\frac{(F_{T}-K)^{+}}{K^{2}}dK
\end{align*}

Plugging this in Equation \ref{eq:replication_with_futures_main},
we have 

\begin{align*}
\frac{1}{T}\int_{0}^{T}\sigma_{t}dt & =\frac{2}{T}(\int_{0}^{F_{0}}\frac{(K-F_{T})^{+}}{K^{2}}dK+\int_{F_{0}}^{\infty}\frac{(F_{T}-K)^{+}}{K^{2}}dK)-\frac{2}{T}\int_{0}^{T}(\frac{1}{F_{0}}-\frac{1}{F_{t}})dF_{t}
\end{align*}

The Black-Scholes assumptions allow us to replicate a variance swap
by holding $\frac{1}{F_{t}}-\frac{1}{F_{0}}$ unites of futures at
$t$ and a set of calls and puts spread across futures price $F_{0}$.
The forward cost of this portfolio is can be written as follows

\begin{align*}
K_{VAR}^{2} & =\frac{2e^{rT}}{T}(\int_{0}^{F_{0}}\frac{P_{0}(K)}{K^{2}}dK+\int_{F_{0}}^{\infty}\frac{C_{0}(K)}{K^{2}}dK)
\end{align*}

Assuming a set of equally spaced strikes $K_{0},...,K_{r},K_{r+1},...,K_{n}$
where $K_{0},K_{1},..,K_{r}$ are strikes for puts and $K_{r+1},K_{r+2},...,K_{n}$
are strikes for calls. Since the options are equally-spaced we have
$\Delta_{K}=K_{i+1}-K_{i}\text{ for all }i\in[1,n-1]$. The replication
portfolio that approximates $\int_{0}^{F_{0}}\frac{P_{0}(K)}{K^{2}}dK+\int_{F_{0}}^{\infty}\frac{C_{0}(K)}{K^{2}}dK$
is 

\begin{align*}
\sum_{i=0}^{r} & \frac{\Delta_{K}P(K_{i})}{K_{i}^{2}}+\sum_{i=r+1}^{n}\frac{\Delta_{K}C(K_{i})}{K_{i}^{2}}
\end{align*}

For a given variance notional exposure $N_{var}$, one must hold $\frac{2\Delta_{K}N_{var}}{TK^{2}}$
of each of option. 

\subsection{\label{subsec:Stocks-based-replication}Stocks-based replication}

The Black-Scholes assumptions readily lead to a replciation algorithm
that is now described (see \cite{demeterfidermanvolswaps1999}). Notice
first that $d(log\ S_{t})=(\mu-\frac{\sigma^{2}}{2})dt+\sigma dZ_{t}$
(for normal variable $Z_{t}$) leads to $\frac{1}{2}\sigma^{2}dt=\frac{dS_{t}}{S_{t}}-d(log\ S_{t})$.
This in turn implies that 

\begin{align}
\frac{1}{T}\int_{0}^{T}\sigma^{2}dt & =\frac{2}{T}(\int_{0}^{T}\frac{dS_{t}}{S_{t}}-log(\frac{S_{T}}{S_{0}}))\label{eq:BlackScholesVariance}
\end{align}

Using $E(\int_{0}^{T}\frac{dS_{t}}{S_{t}})=rT$ and a boundary $S_{*}$
between calls and puts, we make use of the following identity (see
Derman) 

\begin{multline*}
log(\frac{S_{T}}{S_{0}})=log(\frac{S_{*}}{S_{0}})+log(\frac{S_{T}}{S_{*}})\\
=log(\frac{S_{*}}{S_{0}})-(\int_{0}^{S_{*}}\frac{(K-S_{T})^{+}}{K^{2}}+\int_{S_{*}}^{\infty}\frac{(S_{T}-K)^{+}}{K^{2}}-\frac{S_{T}-S_{*}}{S_{*}})
\end{multline*}

The Equation \ref{eq:BlackScholesVariance} becomes 

\begin{align*}
E(\int_{0}^{T}\sigma_{t}^{2}dt) & =\frac{2}{T}(rT-(\frac{S_{0}}{S_{*}}e^{rT}-1)-log(\frac{S_{*}}{S_{0}})+e^{rT}\int_{0}^{S*}\frac{1}{K^{2}}P(K)dK+e^{rT}\int_{S_{*}}^{\infty}\frac{1}{K^{2}}C(K)dK
\end{align*}

Notice that $K_{Var}\equiv E(\int_{0}^{T}\sigma_{t}^{2}dt)$ is called
the strike of the variance swap..

We can rewrite $K_{var}$ as 

\begin{align*}
K_{var} & =\frac{2}{T}(rT-(\frac{S_{0}}{S_{*}}e^{rT}-1)-log(\frac{S_{*}}{S_{0}}))+e^{rT}\Pi_{CP}
\end{align*}
 

The term $\Pi_{CP}$ is the following value of the portfolio of puts
and calls where the weights of the portfolio are calculated using
a recursive algorithm described in Section \ref{subsec:Algorithm-to-replicate-logcontract}

\begin{align}
\Pi_{CP} & =\sum_{i}w(K_{ip})P(S,K_{ip})+\sum_{i}w(K_{ic})C(S,K_{ic})\label{eq:portfolio_derman_cp}
\end{align}


\subsubsection{\label{subsec:Algorithm-to-replicate-logcontract}Algorithm to replicate
Log-contract}

With the portfolio in Equation \ref{eq:portfolio_derman_cp}, let's
consider calls with strikes $K_{1c,}K_{2c}$ etc. such that $S_{*}<K_{1c}<K_{2c}<...$
and puts $K_{1p},K_{2p}$etc. such that $S_{*}>K_{1c}>K_{2c}...$
- where $S^{*}$ is the boundary between calls and puts. We can set
this to initial value of the stock i.e. $S_{*}=S_{0}$.

Let's consider the strike $K_{0}$ at the boundary (reference) stock-price
$S_{*}$. The weight associated with the option struck at $K_{0}$
would be set to 

\begin{align*}
w_{c}(K_{0}) & =\frac{f(K_{1c})-f(K_{0})}{K_{1c}-K_{0}}
\end{align*}

Here, $f(S_{T})$ is defined as

\begin{align*}
f(S_{T}) & =\frac{2}{T}(\frac{S_{T}-S_{*}}{S_{*}}-log(\frac{S_{T}}{S_{*}}))
\end{align*}

The (units) weights for calls in the portfolio are given by 

\begin{align*}
w_{c}(K_{n,c}) & =\frac{f(K_{n+1,c})-f(K_{n,c})}{K_{n,c}-K_{n+1,c}}-\sum_{i=0}^{n-1}w_{c}(K_{i,c})
\end{align*}

Similarly, we can obtain the weights associated with puts by first
setting

\begin{align*}
w_{p}(K_{0}) & =\frac{f(K_{1p})-f(K_{0})}{K_{1p}-K_{0}}
\end{align*}

and then holding other weights of puts determined by

\begin{align*}
w_{p}(K_{n,p}) & =\frac{f(K_{n+1,p})-f(K_{n,p})}{K_{n,p}-K_{n+1,p}}-\sum_{i=0}^{n-1}w_{p}(K_{i,p})
\end{align*}


\subsection{Greeks}

\noindent

The greeks relevant for variance swaps can also be calculated using
the dynamics of Equation \ref{eq:main_replication}.

\bibliographystyle{ieeetr}
\bibliography{fineng}

\end{document}
